%########################################
%CS 461 CAPSTONE 
%GROUP 20 - Griffin Gonsalves, Paul Kwak, Shawn Cross
%FALL 2016
%########################################
\documentclass[letterpaper, 10pt, draftclsnofoot, compsoc, onecolumn]{IEEEtran}
\usepackage[latin1]{inputenc}
\usepackage[T1]{fontenc}
\usepackage[english]{babel}
\usepackage{amsmath}
\usepackage{amssymb,amsfonts,textcomp}
\usepackage{color}
\usepackage{array}
\usepackage{supertabular}
\usepackage{hhline}

\usepackage{hyperref}
%\usepackage{digsig} %not a default package... just for signature field on final page.
\usepackage{comment}

\hypersetup{pdftex, colorlinks=true, linkcolor=black, citecolor=blue, filecolor=blue, urlcolor=blue, pdftitle=SYSTEMS AND SOFTWARE REQUIREMENTS SPECIFICATION (SSRS) TEMPLATE, pdfauthor=Clinton Jeffery, pdfsubject=, pdfkeywords=}
\usepackage[pdftex]{graphicx}
% Outline numbering
\setcounter{secnumdepth}{3}
\renewcommand\thesection{\arabic{section}}
\renewcommand\thesubsection{\arabic{section}.\arabic{subsection}}
\renewcommand\thesubsubsection{\arabic{section}.\arabic{subsection}.\arabic{subsubsection}}

\parindent0pt
\parskip 1.5ex plus 0.2ex minus 0.1ex
\makeatletter

\def\subsubsection{\@startsection{subsubsection}% name
                                 {3}% level
                                 {\z@}% indent (formerly \parindent)
                                 {1ex plus 0.1ex minus 0.1ex}% before skip
                                 {1ex}% after skip
                                 {\normalfont\normalsize}}% style

\newcommand\arraybslash{\let\\\@arraycr}
\makeatother
% Page layout (geometry)
\setlength\voffset{-1in}
\setlength\hoffset{-1in}
\setlength\topmargin{0.5in}
\setlength\oddsidemargin{.75in}
\setlength\evensidemargin{.75in}
\setlength\textheight{8.278in}
\setlength\textwidth{6.5in}
\setlength\footskip{0.561in}
\setlength\headheight{0.5in}
\setlength\headsep{0.461in}
% Footnote rule
\setlength{\skip\footins}{0.0469in}
\renewcommand\footnoterule{\vspace*{-0.0071in}\setlength\leftskip{0pt}\setlength\rightskip{0pt plus 1fil}\noindent\textcolor{black}{\rule{0.25\columnwidth}{0.0071in}}\vspace*{0.0398in}}
% Pages styles
\makeatletter
\newcommand\ps@Standard{
  \renewcommand\@oddhead{\hfill }
  \renewcommand\@evenhead{\@oddhead}
  \renewcommand\@oddfoot{\foreignlanguage{english}{\textcolor{black}{SSRS Page }}\foreignlanguage{english}{\textcolor{black}{\thepage{}}}}
  \renewcommand\@evenfoot{\@oddfoot}
  \renewcommand\thepage{\arabic{page}}
}
\newcommand\ps@FirstPage{
  \renewcommand\@oddhead{}
  \renewcommand\@evenhead{\@oddhead}
  \renewcommand\@oddfoot{}
  \renewcommand\@evenfoot{\@oddfoot}
  \renewcommand\thepage{\arabic{page}}
}
\makeatother
\pagestyle{Standard}
%\setlength\tabcolsep{1mm}
\renewcommand\arraystretch{1.3}
% footnotes configuration
\makeatletter
\renewcommand\thefootnote{\arabic{footnote}}
\makeatother
\title{Forge VR Explorer Requirements}
\author{Shawn Cross, Griffin Gonsalves, Paul Kwak}
\date{2016-11-3}



\begin{document}

\clearpage\setcounter{page}{1}\pagestyle{Standard}
\thispagestyle{FirstPage}

\bigskip

{\centering\selectlanguage{english}\bfseries\color{black}
Technology Review
\par}


\bigskip

{\centering\selectlanguage{english}\bfseries\color{black}
November 10 2016
\par}
\bigskip
\bigskip
\bigskip
\bigskip
\bigskip
\bigskip
\bigskip
\bigskip
\bigskip
\bigskip
\bigskip
\bigskip
%\begin{center}
%	\includegraphics[scale=0.8]{forge_logo.png}
%\end{center}


\vfill
{\centering\selectlanguage{english}\bfseries\color{black}
Abstract
\par}

{\centering\selectlanguage{english}\mdseries\color{black}
	The Forge VR Explorer branches from an Autodesk prototype project called Vrok-It, which is a simple web-based 3D 
	model viewer and mobile virtual reality (VR) explorer. The project will expand upon its ability to display uploaded 3D 
	models in browser and in VR, and improve its accessibility. Conventionally, viewing 3D models in VR is a challenge if 
	you have model files on many devices, or have a headset that only works in conjunction with a smartphone. The 
	Forge VR Explorer aims to do this by utilizing a web-based software that uses the features of the Autodesk Forge API. 
	The project will also be expanded with new ideas and stretch goals as the project is developed.
\par}

{\centering\selectlanguage{english}\bfseries\color{black}
TABLE OF CONTENTS
\par}

\bigskip

{\selectlanguage{english}\bfseries\color{black}
Section\ \ Page}

\setcounter{tocdepth}{9}
\renewcommand\contentsname{}
\tableofcontents

\bigskip
\clearpage


\section[Introduction]{\selectlanguage{english}\rmfamily\bfseries\color{black}
Introduction}

\section[Technologies]{\selectlanguage{english}\rmfamily\bfseries\color{black}
Technologies}

\begin{enumerate}
	\item{File Uploading}
		\begin{enumerate}
			\item{Options}
				\begin{itemize}
					\item  Use a plug-in: Using a plug-in that would handle file uploading would be very convenient as it takes a lot
						of the work out of it. One of the  biggest benefits would be that this could save us a lot of time. If we were to use
						the plug-in we would not have to worry about the effort that would normally go into creating our own, we could simply just
						find one that would fit  what we needed and the put it into our website. However this does also lead to the problem that we
						may not be able to find one that exactly fits what we want. Many of the plug-ins I have found so far simply upload the file to
						a server and then it is done. In our case we need the file to be usable by the website once it is uploaded.
 
					\item  Create our own file upload feature: Creating our own uploading feature would allow us to make the feature 
						work exactly the way that we would want to. Their a many different ways to do this though. One common way I have found 
						is to first create a web form using HTML to make the upload button.Then you could use PHP to submit the file through a POST.
						Using AJAX and JQuery you could also add features that show the progress of the upload and verification such as verifying correct 
						file types are being used. This could take a lot more time and effort to complete though depending on how skilled the person creating 
						the upload feature is and how fancy they want the feature to be.      

					\item  Create our own file uploading feature using new File API added to HTML5: This is similar to the previous option but seems like it might 
						be easier to learn and understand. Many of the features that we wanted are built into the API making it easier for us to create the 
						upload file file feature the way we want it. Though there will be some time for learning how to use the API and some time to completely
						create the feature, there is a lot of good documentation on this API that should make the process go smoothly.  
				\end{itemize}
			\item{Goals}
			\item{Criteria being evaluated}
				\begin{enumerate}
					\item How much time will it take compared to how important it is to the overall project.
					\item How much effort will have to go into learning how to complete the task.
					\item Even if one options isn't exactly what you want does the time benefit outweigh making it exactly the way you wanted.
					\item Which option will offer the most secure way to upload the file.
				\end{enumerate}
			\item{Table}
				\begin{center}
					\begin{tabular}{ | m{10em} | m{10em} | m{10em} | }
						\hline
						Options & Time taken to complete task & Time to learn \\
						
						\hline
						Option 1 & This should take a minimal amount of time to complete this task &  It should also not take a lot of time to  
						learn how to get the plug-in to work with the website \\

						\hline
						Option 2 & This will allow for a more customizable option, however the time to complete the task will be longer & Depending on
						what you are trying to get the feature to do the time to learn could be anywhere from a medium to long period of time  \\
						
						\hline
						Option 3 & This will allow you to create the file upload feature exactly like you want it, it will take a little more time though. 
						 & This shouldn't take a long time to learn since there is a lot of good documentation and examples online \\
						\hline
					\end{tabular}
				\end{center}
			\item{Discussion:} The plug-in option will definitely save the most time in terms of learning it and getting it to work in on the website, 
					however it will likely give you little in the way of customizing and making the file upload feature work exactly how you would want 
					it to. Both of the options to make you own file upload feature will result in a longer time to learn how to do this as well as a longer
					time to complete the task. This does allow you to make sure the the file upload will do exactly what you are wanting it to do. the 
					third option will is also using HTML5 this means that is will work in all the new browsers.
			\item{Select Option:}
					 I believe that the extra time that it takes to create our own custom file upload feature will outweigh the extra amount of time 
					that it will take to do this. I also think that the use of the file API is better since it is well documented, easy to find examples
					for, and has many of the features that we wanted built into it.
		\end{enumerate}
	\item{FIle conversion}
		\begin{enumerate}
			\item{Option}
				\begin{itemize}
					\item accusoft API: This is an API that allows users to convert their CAD files in to raster files such as svg (scalable vector graphic) files that could 
						to view in the model viewer. This API does cost money and depending on how many conversions the user is planning on doing this option
						could turn out to be relatively expensive. There is not a lot of documentation on how this API can be integrated into a application so 
						that may make it take longer to get the task complete.

					\item  Forge Model derivative API: This API has the ability to take users CAD source file and convert them into OBJ and STL files. This API can also
						convert the files directly into SVF files the same type of file that the model viewer API we will be using on this project uses. This API would 
						also not cost us anything to use while we are working on the project since the company that created the API is the same company we will 
						be making the project for. There might be a little bit of a learning curve when it comes to integrating it into the website but the API has a 
						lot of tutorials and documentation on how to use the API in an application.  This API also uses a token based authentication system that
						provides security when converting the users files. 

					\item  cloudconvert API: This API gives you the ability to convert from CAD source files to a SVG file but it also give you the ability to convert the 
						other way was well. This API also has well documented instruction on how to integrate this API into an application. That would make it easier
						to learn and take less time to complete the task. This API would cost money to use depending on how much time is spent doing the conversions.
			
				\end{itemize}
			\item{Goals}
			\item{Criteria being evaluated}
				\begin{enumerate}
					\item How much time will it take compared to how important it is to the overall project.
					\item How much effort will have to go into learning how to complete the task.
					\item How secure is the file conversion.
				\end{enumerate}
			\item{Table?}
			\item{Discussion:} These three options all offer similar conversions. They all convert the CAD source files into a usable file for the model viewer. The Forge API will however 
				convert the file into an SVF file that can then be directly loaded into the viewer. This could save time and make it easier to complete the task. The Forge API is also 
				free for us to use during the duration of the project so that would also make a big difference as well. Finally the Forge API offers the most secure option when it comes
				to converting the users files since it  requires a token authentication to even do the conversion. This token is created as soon as the user visits the website and last for 
				the duration that the user is on the website. Also with lots of documentation as well as tutorials on how to integrate the API into our application I feel like this option will
				take the least amount of time to complete the task. 
			\item{Select Option:}
				I Think that we should use the Forge model derivative API since it seems the most straight forward on how to integrate it into our application. It won't cost us any money
				to use and has a more secure way of converting the users files. 
		\end{enumerate}
	\item{Technology 3}
		\begin{enumerate}
			\item{Option}
			\item{Goals}
			\item{Criteria being evaluated}
			\item{Table?}
			\item{Discussion}
			\item{Select Option}
		\end{enumerate}
	\item{Technology 4}
		\begin{enumerate}
			\item{Option}
			\item{Goals}
			\item{Criteria being evaluated}
			\item{Table?}
			\item{Discussion}
			\item{Select Option}
		\end{enumerate}
	\item{Technology 5}
		\begin{enumerate}
			\item{Option}
			\item{Goals}
			\item{Criteria being evaluated}
			\item{Table?}
			\item{Discussion}
			\item{Select Option}
		\end{enumerate}
	\item{Technology 6}
		\begin{enumerate}
			\item{Option}
			\item{Goals}
			\item{Criteria being evaluated}
			\item{Table?}
			\item{Discussion}
			\item{Select Option}
		\end{enumerate}
	\item{Technology 7}
		\begin{enumerate}
			\item{Option}
			\item{Goals}
			\item{Criteria being evaluated}
			\item{Table?}
			\item{Discussion}
			\item{Select Option}
		\end{enumerate}
	\item{Technology 8}
		\begin{enumerate}
			\item{Option}
			\item{Goals}
			\item{Criteria being evaluated}
			\item{Table?}
			\item{Discussion}
			\item{Select Option}
		\end{enumerate}
	\item{Technology 9}
		\begin{enumerate}
			\item{Option}
			\item{Goals}
			\item{Criteria being evaluated}
			\item{Table?}
			\item{Discussion}
			\item{Select Option}
		\end{enumerate}
\end{enumerate}

\section[Conclusion]{\selectlanguage{english}\rmfamily\bfseries\color{black}
Conclusion}

\section[Bibliography]{\selectlanguage{english}\rmfamily\bfseries\color{black}
Bibliography}
\end{document}
